% ============================================================
% audit_table.tex -- Paper III: Tidal Ionization / Quantum Roche Limit
%
% A "system audit" table mapping the paper's claims to the evidence
% backing each one.  This is *the* canonical source of truth for which
% claims are PASS / PART / STUB / FAIL --- prose elsewhere in the paper
% should not contradict it.
%
% Status semantics:
%   PASS  -- derivation complete or experiment confirms the claim to
%            stated precision.
%   PART  -- mechanism demonstrated, full quantitative match pending;
%            specific gaps documented in the evidence column.
%   STUB  -- analytical result with no numerical confirmation yet,
%            or planned experiment / paper-text derivation not yet
%            written.
%   FAIL  -- tested and disconfirmed.
% ============================================================

\label{tab:audit}

{\scriptsize
\setlength{\tabcolsep}{4pt}
\renewcommand{\arraystretch}{0.95}
\begin{longtable}{@{}p{2.8cm}p{3.6cm}p{3.6cm}p{4.0cm}p{1.0cm}@{}}
\caption{\small Audit for Paper~III.  Inherited rows are evidenced by
Paper~I material (analytical derivations in
\texttt{notes/plasma\_as\_gravitational\_ionization.md}; numerical
machinery in \texttt{exp\_18\_tidal\_ionization.py}, reproduced
verbatim in this repo).  Paper~III-specific rows -- the explicit
$M_\text{min}(d)$ derivation, the $d^3$ scaling identification, the
g-scan run results, and the observational predictions -- start
\texttt{STUB}/\texttt{PART} in v0.1 and flip as Phase~2 lands.}
\label{tab:audit_caption} \\
\hline
\textbf{Claim} & \textbf{Mechanism} & \textbf{Continuum / formal limit} & \textbf{Evidence} & \textbf{Status} \\
\hline
\endfirsthead

\multicolumn{5}{l}{\textit{(Table~\ref{tab:audit}, continued from previous page)}} \\
\hline
\textbf{Claim} & \textbf{Mechanism} & \textbf{Continuum / formal limit} & \textbf{Evidence} & \textbf{Status} \\
\hline
\endhead

\hline
\multicolumn{5}{r}{\textit{(continued on next page)}} \\
\endfoot

\hline
\endlastfoot

% ── Inherited from Paper I ──────────────────────────────────────────
\multicolumn{5}{l}{\textit{Inherited from Paper~I~\cite{menendez2026geometry}.}} \\
\hline

Hydrogen as joint two-session Arnold-tongue lock-in
    & Bipartite tick rule + clock-density mass; lock-in at
      $\omega \cdot R_1 = \pi/3$
    & Standard hydrogen spectrum in the continuum limit
    & \texttt{exp\_12} (Paper~I), $k_\text{min} = 0.0970$ vs
      $k_\text{Bohr} = 0.0971$ to 4 sig figs
    & \texttt{PASS} \\

Gradient detuning breaks the joint resonance
    & $\omega_\text{eff}(x) = \omega \cdot \bar{\rho} /
      \rho_\text{clock}(x)$; differential detuning across $R_1$
      exceeds $\Delta\omega_\text{tongue}$
    & Three ionization channels (thermal Saha + pressure-degenerate
      + gradient detuning); the third is new in this framework
    & \texttt{notes/plasma\_as\_gravitational\_ionization.md}
      (analytical)
    & \texttt{PASS} \\

% ── Paper III-specific ──────────────────────────────────────────────
\hline
\multicolumn{5}{l}{\textit{Paper~III-specific claims (this paper).}} \\
\hline

Quantum Roche limit $M_\text{min}(d) \cdot R_1 / d^3 =
\Delta\omega_\text{tongue}$
    & Substitute $g = M / d^2$ into the gradient-detuning condition;
      solve for $M$ at threshold
    & $M_\text{min}(d) = \Delta\omega_\text{tongue} \cdot d^3 / R_1$;
      same $d^3$ scaling as classical Roche limit~\cite{roche1849}
      with $\Delta\omega_\text{tongue}$ replacing body density
    & Paper-text derivation pending (\S\,planned)
    & \texttt{STUB} \\

$d^3$ scaling is the quantum-mechanical analogue of the classical
Roche limit
    & Both derive from $|\nabla F| \cdot \ell / F = $ const at
      threshold, where $\ell$ is the bound-state size and $F$ is
      the strength of the binding force; the framework derives the
      quantum case from the Arnold-tongue width
    & The two limits are the same identity seen from two levels of
      description; in the lattice they are the same thing
    & Paper-text observation pending (\S\,planned)
    & \texttt{STUB} \\

Three ionization channels (thermal, pressure, gradient)
    & Saha (kinetic) + Pauli (degenerate) + Arnold-tongue
      detuning (gradient); the third is new in the framework
    & Standard astrophysics recognises the first two; the third
      becomes dominant near compact objects
    & \texttt{notes/plasma\_as\_gravitational\_ionization.md};
      paper-text consolidation pending
    & \texttt{PART} \\

Numerical confirmation:
$T_\text{escape}(g)$ decreases monotonically with $g$, sharp drop
near $g_\text{crit}$
    & g-scan in \texttt{exp\_18}: apply linear gradient
      $V_\text{tidal} = g \cdot (x - x_\text{centre})$ from
      $t = 0$; measure tick at which $r_\text{pdf} > 2 R_1$ for
      $N$ consecutive windows
    & $g_\text{crit} \approx \Delta\omega_\text{tongue} / R_1$;
      Paper~I's free-particle tongue width gives
      $g_\text{crit} \approx 0.016$ per node
    & \texttt{exp\_18\_tidal\_ionization.py}: script in place
      (Paper~I copy, reproduced here); runtime evidence pending
      first run from this repo
    & \texttt{STUB} \\

Tidal displacement diagnostic:
electron and proton move in opposite directions under gradient
(not isotropically as for thermal disruption)
    & x-projection of CoM relative to system CoM should show
      $x_e \cdot x_p < 0$ for gradient ionization, isotropic
      scatter for energy-injection disruption
    & Distinguishes gradient ionization from
      kinetic-energy-injection mechanisms
    & \texttt{exp\_18} diagnostic columns
      (\texttt{x\_e\_proj}, \texttt{x\_p\_proj}); runtime evidence
      pending
    & \texttt{STUB} \\

Atomic ISCO: innermost stable atomic orbit
$r_\text{atomic\_ISCO}(M)$
    & Largest $d$ at which gradient ionization always succeeds
      for the $n=1$ ground state, independent of temperature;
      analogue of geodesic ISCO at $r = 6 GM / c^2$
    & Innermost neutral-hydrogen radius around compact objects;
      structural prediction independent of thermal modelling
    & Paper-text derivation pending (\S\,planned)
    & \texttt{STUB} \\

Observational signature:
21\,cm HI inner-edge morphology near compact objects
    & Threshold $M_\text{min}(d)$ gives a temperature-independent
      inner boundary for neutral hydrogen
    & 21\,cm HI absorption / emission inner edge near neutron stars
      and active galactic nuclei
    & Paper-text comparison to existing 21\,cm HI surveys pending
      (\S\,planned)
    & \texttt{STUB} \\

Observational signature:
Fe K$\alpha$ line-profile inner edge in accretion disks
    & Atomic ISCO gives an inner edge for Fe atoms supporting
      K$\alpha$ emission; below the atomic ISCO, Fe is fully
      ionized regardless of temperature
    & Inner edge of Fe K$\alpha$ emission near accretion disks
      (XMM-Newton / NICER profiles)
    & Paper-text comparison pending (\S\,planned)
    & \texttt{STUB} \\

\hline
\end{longtable}
}
